\chapter*{LỜI MỞ ĐẦU}
\addcontentsline{toc}{chapter}{LỜI MỞ ĐẦU}

\section*{1. Lý do chọn đề tài}

Trong bối thương mại điện tử tại Việt Nam phát triển mạnh mẽ, đặc biệt trong lĩnh vực thời trang trực tuyến. Tuy nhiên, một trong những hạn chế lớn nhất của việc mua sắm quần áo qua mạng là người dùng không thể thử sản phẩm trước khi đặt hàng, dẫn đến tỷ lệ đổi trả cao và làm giảm trải nghiệm mua sắm. Bên cạnh đó, nhiều cửa hàng thời trang quy mô vừa và nhỏ vẫn phụ thuộc vào các sàn thương mại điện tử lớn, gây khó khăn trong việc xây dựng thương hiệu riêng và quản lý dữ liệu khách hàng.

Sự phát triển của AI, Computer Vision và Augmented Reality trong những năm gần đây đã mở ra khả năng xây dựng các hệ thống thử đồ ảo với chi phí triển khai phù hợp. Các mô hình AI Virtual Try-On hiện đại như FitDiT, kết hợp với các kỹ thuật phát hiện keypoint cơ thể và Human Parsing, cho phép sinh ảnh thử đồ với chất lượng cao. Đồng thời, công nghệ AR hỗ trợ hiển thị quần áo trực tiếp trên cơ thể người dùng thông qua camera theo thời gian thực, góp phần nâng cao tính trực quan trong quá trình mua sắm.

Xuất phát từ thực tế đó, đề tài lựa chọn xây dựng một hệ thống thương mại điện tử thời trang tích hợp AI Virtual Try-On, AR Try-On và chatbot AI nhằm hỗ trợ tư vấn sản phẩm, nâng cao trải nghiệm người dùng và hỗ trợ hoạt động quản lý kinh doanh cho cửa hàng thời trang.

\section*{2. Tình hình nghiên cứu}

Trên thế giới, nhiều nền tảng thương mại điện tử lớn như Amazon, Zara và ASOS đã bước đầu tích hợp công nghệ Virtual Try-On vào quy trình mua sắm trực tuyến nhằm nâng cao trải nghiệm khách hàng. Trong lĩnh vực nghiên cứu học thuật, các mô hình như VITON, HR-VITON và FitDiT dựa trên kiến trúc diffusion đã cho thấy hiệu quả cao trong bài toán sinh ảnh thử đồ ảo.

Song song đó, các kỹ thuật phát hiện keypoint cơ thể như DWpose và Human Parsing hỗ trợ quá trình phân tích hình dạng cơ thể và phân vùng quần áo chính xác hơn, giúp cải thiện chất lượng kết quả thử đồ. Trong lĩnh vực AR, các nền tảng như Snap Camera Kit hỗ trợ triển khai thử đồ trực tiếp trên trình duyệt thông qua camera thời gian thực.

Ngoài ra, các mô hình ngôn ngữ lớn (LLM) và kỹ thuật Retrieval-Augmented Generation (RAG) cũng đang được ứng dụng rộng rãi trong các hệ thống chatbot hỗ trợ khách hàng. Việc kết hợp vector database với mô hình ngôn ngữ giúp chatbot có khả năng truy xuất dữ liệu ngữ nghĩa và phản hồi phù hợp với ngữ cảnh người dùng.

Tại Việt Nam, các sàn thương mại điện tử lớn như Shopee hay Lazada chủ yếu ứng dụng AI vào gợi ý sản phẩm và tìm kiếm thông minh, trong khi các hệ thống tích hợp đồng thời AI Virtual Try-On, AR Try-On và chatbot AI trong cùng một nền tảng thương mại điện tử vẫn chưa được khai thác phổ biến, đặc biệt đối với các cửa hàng thời trang quy mô vừa và nhỏ.

\section*{3. Mục đích nghiên cứu}

Đề tài hướng đến việc xây dựng một hệ thống thương mại điện tử thời trang tích hợp các công nghệ AI và AR nhằm nâng cao trải nghiệm mua sắm trực tuyến cho người dùng. Hệ thống hỗ trợ thử quần áo ảo bằng AI Virtual Try-On dựa trên mô hình FitDiT kết hợp DWpose và Human Parsing, đồng thời hỗ trợ thử đồ AR theo thời gian thực thông qua Snap Camera Kit.

Ngoài ra, hệ thống còn tích hợp chatbot AI sử dụng mô hình ngôn ngữ lớn chạy cục bộ thông qua Ollama kết hợp kỹ thuật RAG và PostgreSQL pgvector nhằm hỗ trợ tư vấn sản phẩm và tương tác với khách hàng. Hệ thống đồng thời cung cấp các chức năng quản lý sản phẩm, đơn hàng, thanh toán và quản trị dành cho cửa hàng thời trang.

\section*{4. Nhiệm vụ nghiên cứu}

\begin{itemize}
    \item Phân tích yêu cầu và thiết kế kiến trúc tổng thể cho hệ thống thương mại điện tử thời trang tích hợp AI và AR.
    
    \item Xây dựng giao diện người dùng bằng Vue.js, Nuxt.js và TailwindCSS, bảo đảm khả năng responsive trên nhiều thiết bị khác nhau.
    
    \item Phát triển backend bằng FastAPI theo kiến trúc RESTful API, thiết kế cơ sở dữ liệu quan hệ trên PostgreSQL và sử dụng pgvector để lưu trữ vector embedding.
    
    \item Xây dựng chức năng AI Virtual Try-On sử dụng mô hình FitDiT kết hợp DWpose và Human Parsing nhằm sinh ảnh thử đồ ảo cho người dùng.
    
    \item Xây dựng chức năng AR Try-On sử dụng Snap Camera Kit và mô hình 3D để hiển thị quần áo trực tiếp trên cơ thể người dùng thông qua camera.
    
    \item Xây dựng chatbot AI sử dụng Ollama kết hợp kỹ thuật RAG nhằm hỗ trợ tư vấn sản phẩm và trả lời câu hỏi của khách hàng.
    
    \item Tích hợp thanh toán trực tuyến bằng Stripe và hỗ trợ thanh toán khi nhận hàng.
    
\end{itemize}

\section*{5. Phương pháp nghiên cứu}

\begin{itemize}
    \item Thu thập, đọc và phân tích các tài liệu kỹ thuật, documentation chính thức của các framework, thư viện và mô hình được sử dụng như Vue.js, Nuxt.js, FastAPI, FitDiT, DWpose, Ollama, PostgreSQL pgvector và Snap Camera Kit.
    
    \item Áp dụng phương pháp phân tích và thiết kế hệ thống theo hướng module hóa, tách biệt frontend, backend và AI service nhằm tăng khả năng mở rộng và bảo trì hệ thống.
    
    \item Áp dụng Firebase Authentication kết hợp JWT, RBAC và cơ chế Rate Limiting trong FastAPI nhằm kiểm soát truy cập và hạn chế các yêu cầu bất thường đến hệ thống.    \item Sử dụng các mô hình học sâu đã được huấn luyện sẵn như FitDiT, DWpose và Human Parsing thay vì huấn luyện từ đầu nhằm tập trung vào bài toán tích hợp hệ thống và tối ưu trải nghiệm người dùng.
    
    \item Áp dụng kỹ thuật RAG cho chatbot AI bằng cách tạo embedding dữ liệu sản phẩm, lưu trữ vector trên PostgreSQL pgvector và sinh phản hồi bằng mô hình ngôn ngữ chạy cục bộ thông qua Ollama.
    
    \item Kiểm thử API bằng Postman, kiểm thử giao diện trên nhiều trình duyệt như Chrome, Firefox và Edge, đồng thời kiểm thử trên cả desktop và mobile nhằm đánh giá khả năng vận hành thực tế của hệ thống.
    
    \item Sử dụng Git để quản lý mã nguồn.
\end{itemize}

\section*{6. Kết quả đạt được}

\begin{itemize}
    \item Xây dựng hoàn chỉnh hệ thống thương mại điện tử thời trang với đầy đủ các chức năng quản lý sản phẩm, giỏ hàng, đặt hàng, thanh toán và phân quyền người dùng.
    
    \item Tích hợp chức năng AI Virtual Try-On sử dụng FitDiT, DWpose và Human Parsing, cho phép người dùng thử quần áo ảo thông qua hình ảnh cá nhân.
    
    \item Xây dựng chức năng AR Try-On hỗ trợ hiển thị quần áo trên cơ thể người dùng thông qua camera theo thời gian thực.
    
    \item Xây dựng chatbot AI sử dụng Ollama kết hợp kỹ thuật RAG và PostgreSQL pgvector nhằm hỗ trợ tư vấn và gợi ý sản phẩm.
    
    \item Tích hợp thanh toán trực tuyến bằng Stripe và hỗ trợ thanh toán COD.
    
    \item Hệ thống hỗ trợ xác thực JWT, phân quyền người dùng theo vai trò và giới hạn tần suất truy cập API nhằm nâng cao tính bảo mật.
    
    \item Hệ thống được kiểm thử trên nhiều trình duyệt và thiết bị khác nhau, bảo đảm khả năng vận hành ổn định trong môi trường thực tế.
\end{itemize}

\section*{7. Kết cấu luận văn tốt nghiệp}

Luận văn tốt nghiệp được tổ chức thành 4 chương như sau:

\begin{itemize}
    \item Chương 1: Cơ sở lý thuyết
    \item Chương 2: Phân tích, thiết kế và xây dựng hệ thống
    \item Chương 3: Cài đặt và kiểm thử
    \item Chương 4: Kết luận và hướng phát triển
\end{itemize}


\clearpage